% Brief (tikz-diagram-craft).
% Reader question: how long does a marker live, and how sensitive is that to r?
% Claim: moving the decay rate from 0.9 to 0.95 more than doubles a marker's
%   life (pruned at tick 29 vs tick 59) against the 0.05 prune threshold.
% Evidence: w(t) = r^t, threshold 0.05; 0.9^28 = .052 > .05 > .047 = 0.9^29;
%   0.95^58 = .051 > .05 > .0485 = 0.95^59; ln .05 / ln .95 = 58.4. [analytic]
% Must distinguish: the two rates; where each crosses the threshold; lifetime.
% Grammar: xy plot with the crossing marked and dropped to the tick axis, the two
%   lifetimes as brackets under the axis. Rejected: a table of (r, lifetime),
%   which states the numbers but hides that the curve is flat near the
%   threshold -- the reason small changes in r move the crossing so far.
% Five-second test: "the flatter curve crosses the dashed line twice as late".
\begin{figure}[H]
\centering
\pdfigurehue{pdcobalt}%
\begin{tikzpicture}[pd figure]
\begin{axis}[pd axis,
  width=7.7cm,height=4.4cm,
  xmin=0,xmax=64,ymin=0,ymax=1,
  xtick={0,10,20,29,40,50,59},ytick={0,0.25,0.5,0.75,1},
  yticklabels={0,.25,.5,.75,1},
  xlabel={ticks since the marker was deposited},
  ylabel={marker weight $w(t)$}]
  % the threshold: a warning rule, labelled in ink
  \addplot[pd warn series,line width=.9pt,domain=0:64] {0.05};
  \node[pd label,anchor=west,align=left] at (axis cs:64,0.05) {prune at\\ $w<0.05$};
  % the two rates
  \addplot[pd series,mark=none,domain=0:64,samples=129] {0.9^x};
  \addplot[pd focus series,mark=none,domain=0:64,samples=129] {0.95^x};
  \node[pd label,anchor=north east] at (axis cs:11,{0.9^11}) {$r=0.9$};
  \node[pd focus label,anchor=south west] at (axis cs:17,{0.95^17}) {$r=0.95$};
  % the crossings, dropped to their ticks
  \draw[pd guide] (axis cs:29,0.047) -- (axis cs:29,0);
  \draw[pd guide] (axis cs:59,0.0485) -- (axis cs:59,0);
  \node[pd datum] at (axis cs:29,0.047) {};
  \node[pd focus datum] at (axis cs:59,0.0485) {};
\end{axis}
\end{tikzpicture}
\caption{With a prune threshold of $0.05$, raising decay retention from $0.9$ to $0.95$ delays pruning from tick 29 to tick 59. [internal; analytic]}
\label{fig:swk-marker-decay}
\end{figure}
